\chapter{CP/M: the Z80 Second Processor}\label{cha:cpm}

In 1984 Acorn sold a Z80 Second Processor for the BBC Micro; its purpose
was to run CP/M, the operating system that owned business computing before
the IBM PC. It sat on the Tube. So does ours. \verb|CPM| at the shell boots
CP/M~2.2 on a Z80 co-processor (RunCPM, on the host), and the \verb|A0>|
prompt that defined a decade appears. The boot banner cheerfully measures
the Z80 in \emph{gigahertz} --- read that number knowingly: it is the
emulated Z80's speed \emph{on whatever host you run the emulator on} (a
laptop, in this book's screenshots), so it varies machine to machine.
For scale, the real thing shipped at 4\,MHz.

\begin{type}
CPM
\end{type}

\screen{shots/cpm.png}{CP/M 2.2 over the Tube: the CCP's \texttt{DIR} and
\texttt{TYPE}, on drive A.}

\section{Drives are folders}
Drives \verb|A:| to \verb|P:| are the host folders \texttt{fs/CPM/A}
through \texttt{fs/CPM/P}; CP/M's user areas 0--15 are numbered subfolders.
Drop any \texttt{.COM} file into \texttt{fs/CPM/A/0} and run it by name.
\verb|DIR|, \verb|TYPE|, \verb|ERA|, \verb|REN|, \verb|SAVE| and
\verb|USER| are built into the CCP; \verb|EXIT| hands the console back to
the K/OS shell.

\section{Software}
The RunCPM project ships a complete system disk
(\texttt{DISK/A0.zip} in its repository): Digital Research's own
\texttt{ASM}, \texttt{MAC}, \texttt{DDT}, \texttt{ZSID}, \texttt{STAT},
\texttt{PIP} and \texttt{ED}, the \texttt{SUBMIT} batch system, Microsoft
BASIC, a Z80 assembler with its manual, XMODEM, and the source code of the
BDOS and CCP --- buildable on the machine itself. One
\texttt{unzip A0.zip -d fs/CPM/} installs the lot. Beyond that lies the
whole surviving CP/M software world: WordStar, Turbo Pascal, dBase~II,
Zork, one archive away.

\begin{aside}
\textbf{Edges, honestly:} line-mode programs (assemblers, compilers,
MBASIC, adventures) work today, and so does the full-screen software:
the console is JIM, a VT100 in hardware (Chapter~\ref{cha:tube}).
WordStar~4 on \texttt{E:} user~3 comes installed for ``ANSI Standard''
at 79$\times$29 --- \verb|WS| and you are editing; Turbo Pascal~3 on
\texttt{H:} user~3 needs nothing either. A program that asks its
terminal type at install time wants ``VT100'' or ``ANSI''.

\screen{shots/wordstar.png}{WordStar~4 editing its own \texttt{READ.ME}:
the edit menu, the ruler, the text --- 79$\times$29 on JIM.}
Desktop host only, like the rest of the Tube.
\end{aside}
